\documentclass[11pt,a4paper]{article}

% ============================================================
% Paper 4-2 (versi Bahasa Indonesia) --- draft kerja.
% Judul: Robust Aggregation untuk Federated NIDS lintas-dataset melawan
%        poisoning & backdoor, di atas ruang SFM (fondasi Paper 1, infra Paper 4).
% CATATAN KEJUJURAN: seluruh ANGKA HASIL masih [TBD] -- diisi dari eksperimen NYATA
%   (notebook notebook42/42_01..42_07 + eksekusi AWS). JANGAN mengarang angka.
% ============================================================

\usepackage[utf8]{inputenc}
\usepackage[T1]{fontenc}
\usepackage[bahasa]{babel}
\usepackage{graphicx}
\usepackage{algorithm}
\usepackage{algpseudocode}
\usepackage{booktabs}
\usepackage{amsmath,amssymb}
\usepackage{siunitx}
\usepackage{hyperref}
\usepackage{geometry}
\usepackage{caption}
\usepackage{multirow}
\usepackage{placeins}
\geometry{margin=2.5cm}
\graphicspath{{figure-fed42/}}
\hypersetup{colorlinks=true, linkcolor=blue, citecolor=blue, urlcolor=blue}

\newcommand{\TBD}{\textit{[TBD]}}

\title{\textbf{Agregasi Tangguh untuk Sistem Deteksi Intrusi Jaringan Terfederasi
Lintas-\emph{Dataset}: Pertahanan terhadap Serangan \emph{Poisoning} dan
\emph{Backdoor} pada Ruang Fitur Semantik}}

\author{Hero Yudo Martono\thanks{Penulis koresponden: \texttt{hero@pens.ac.id}},
\quad Iwan Syarif, \quad Ferry Astika Saputra \\[0.3em]
\textit{Departemen Teknik Informatika, Politeknik Elektronika Negeri Surabaya,
Surabaya, Indonesia} \\
\texttt{hero@pens.ac.id; iwanarif@pens.ac.id; ferryas@pens.ac.id}}

\date{Draft --- \today}

\begin{document}
\maketitle

\begin{abstract}
\emph{Federated learning} (FL) memungkinkan banyak organisasi melatih sistem
deteksi intrusi jaringan (NIDS) secara kolaboratif tanpa memusatkan data trafik
sensitif. Namun agregasi standar FedAvg mengasumsikan seluruh klien jujur ---
asumsi yang berbahaya pada keamanan siber, karena sebagian klien dapat
dikompromikan dan mengirim pembaruan model beracun (\emph{poisoning}) atau menanam
\emph{backdoor} (pola pemicu yang memaksa label target). Makalah ini mengkaji
\emph{robust aggregation} untuk FL-NIDS pada setting yang jarang diuji:
\textbf{lintas-\emph{dataset}} (CSE-CIC-IDS2018 dan UNSW-NB15) yang disatukan dalam
ruang \emph{Semantic Feature Mapping} (SFM) dari karya pendahulu kami \cite{sfm}.
Kami membandingkan FedAvg terhadap empat aturan tangguh --- Krum, Multi-Krum,
\emph{Trimmed-Mean}, dan \emph{Median} per-koordinat --- di bawah tiga model
ancaman (\emph{label-flipping}, penskalaan bobot, dan \emph{backdoor}) dengan
fraksi penyerang 0--30\%. Evaluasi memakai \emph{Matthews Correlation Coefficient}
(MCC) untuk kualitas deteksi dan \emph{Attack Success Rate} (ASR) untuk efektivitas
\emph{backdoor}, dengan protokol multi-\emph{seed}. \textbf{Catatan:} seluruh angka
hasil pada draf ini masih \TBD{} karena eksperimen belum dieksekusi; nilai akan
diisi dari hasil nyata sebelum diajukan untuk publikasi.

\medskip
\noindent\textbf{Kata kunci:} deteksi intrusi jaringan; \emph{federated learning};
\emph{robust aggregation}; \emph{poisoning}; \emph{backdoor}; keamanan model.
\end{abstract}

% ============================================================
\section{Pendahuluan}
\label{sec:intro}

\emph{Federated learning} telah menjadi arah penting untuk NIDS kolaboratif karena
memungkinkan pelatihan model bersama tanpa memindahkan data mentah lintas-organisasi
\cite{popoola,fsurvey1}. Namun kemampuan berbagi model ini membuka permukaan
serangan baru: dalam FL, setiap klien menyumbang pembaruan parameter, dan klien
yang dikompromikan dapat mengirim pembaruan yang dimanipulasi. Jika satu atau
beberapa klien berperilaku jahat, model global dapat terdegradasi, atau --- lebih
berbahaya --- tertanam \emph{backdoor} yang membuat trafik serangan tertentu lolos
sebagai sah \cite{fbackdoor,fbagdasaryan}.

FedAvg \cite{popoola}, sebagai agregasi baku, hanya merata-ratakan pembaruan
berbobot ukuran data; ia tidak memiliki mekanisme untuk membedakan kontribusi jujur
dari yang beracun. Pada NIDS, kerentanan ini kritis: data jaringan sangat
heterogen dan tersebar, sehingga pembaruan yang menyimpang sulit dibedakan dari
variasi wajar antar-klien. Kebutuhan akan \emph{robust aggregation} --- aturan
agregasi yang menahan pengaruh klien jahat --- menjadi nyata.

Makalah ini mengangkat satu celah spesifik. Mayoritas kajian \emph{robust
aggregation} \cite{fkrum,ftrimmed} dievaluasi pada data IID atau \emph{benchmark}
umum. Sementara itu, karya pendahulu kami pada Paper~4 menunjukkan bahwa FL-NIDS
\textbf{lintas-\emph{dataset}} (klien dengan alat ekstraksi fitur berbeda, disatukan
via SFM) sudah menantang bahkan \emph{tanpa} penyerang. Pertanyaan yang belum
terjawab: bagaimana interaksi antara heterogenitas lintas-\emph{dataset} ini dengan
kehadiran klien jahat, dan apakah \emph{robust aggregation} tetap efektif pada
setting ini?

\paragraph{Pertanyaan riset.}
\begin{itemize}
  \item \textbf{RQ1.} Seberapa rapuh FedAvg pada FL-NIDS lintas-\emph{dataset}
  (ruang SFM) seiring naiknya fraksi penyerang?
  \item \textbf{RQ2.} Sejauh mana \emph{robust aggregation} (Krum, Multi-Krum,
  \emph{Trimmed-Mean}, \emph{Median}) memulihkan MCC dan menekan ASR \emph{backdoor}?
  \item \textbf{RQ3.} Berapa ``harga ketahanan'' --- penurunan MCC \emph{robust
  aggregation} ketika tidak ada serangan?
\end{itemize}

\paragraph{Kontribusi.}
\begin{enumerate}
  \item Evaluasi sistematis kerentanan FedAvg dan empat aturan \emph{robust
  aggregation} pada FL-NIDS \textbf{lintas-\emph{dataset} di ruang SFM} --- setting
  yang lebih realistis dan lebih sulit daripada \emph{benchmark} IID umum.
  \item Tiga model ancaman konkret (\emph{label-flipping}, penskalaan bobot,
  \emph{backdoor}) dengan metrik ganda: MCC untuk deteksi dan ASR untuk \emph{backdoor}.
  \item Analisis jujur \emph{trade-off} ketahanan versus kualitas deteksi
  (RQ3), dengan protokol multi-\emph{seed}.
\end{enumerate}

\paragraph{Posisi dalam silsilah riset kami.} Paper~1 memperkenalkan SFM
(generalisasi lintas-jaringan, terpusat); Paper~2 ketahanan adversarial (terpusat);
Paper~3 adaptasi \emph{online}; Paper~4 memindahkan SFM ke \emph{federated} dan
mengkaji interoperabilitas skema-fitur. Makalah ini (Paper~4-2) membuka sumbu
\textbf{keamanan/ketahanan} dalam federasi: bukan lagi ``dapatkah klien berskema
beda berfederasi'', melainkan ``dapatkah federasi itu bertahan ketika sebagian
klien jahat''. Paper~4 sendiri menyebut pertahanan \emph{poisoning} sebagai kerja
mendatang; makalah ini mengisinya.

% ============================================================
\FloatBarrier
\section{Karya Terkait}
\label{sec:related}

\paragraph{FL untuk NIDS.} FL-NIDS dikaji luas dari sisi privasi, efisiensi
komunikasi, dan heterogenitas \cite{fsurvey1,fsurvey2,ftaxonomy}. Sebagian besar
mengasumsikan klien tepercaya, sehingga aspek keamanan agregasi kurang tergarap.

\paragraph{Serangan terhadap FL.} \emph{Poisoning} dan \emph{backdoor} pada FL
telah ditunjukkan efektif \cite{fbackdoor,fbagdasaryan,ffung}: sedikit klien jahat
dapat menggeser model global atau menanam pemicu tersembunyi. Pada NIDS, dampaknya
berupa serangan yang lolos deteksi atau lonjakan \emph{false alarm}.

\paragraph{Robust aggregation.} Aturan tangguh seperti Krum \cite{fkrum},
\emph{Trimmed-Mean} dan \emph{Median} \cite{ftrimmed}, serta skema berbasis
kepercayaan (\emph{trust}) ditawarkan untuk meredam \emph{outlier}. Namun sebagian
besar dievaluasi pada data IID/\emph{benchmark} umum. Penerapan dan evaluasi
spesifik pada FL-NIDS \textbf{lintas-\emph{dataset}} (klien berskema beda via SFM)
masih tipis --- inilah ruang kontribusi kami.

\paragraph{Set fitur bersama.} Upaya menuju ruang fitur NIDS bersama
\cite{sarhan} dan SFM \cite{sfm} memungkinkan federasi lintas-\emph{tool}; makalah
ini memeriksa keamanan federasi di ruang bersama tersebut.

% ============================================================
\FloatBarrier
\section{Metode}
\label{sec:method}

\subsection{Basis: FL-NIDS di ruang SFM}
\label{sec:base}
Model lokal adalah \emph{multilayer perceptron} ringan $9\!\rightarrow\![128,64]\!\rightarrow\!1$
pada sembilan fitur SFM, dengan normalisasi \emph{z-score} per-klien lokal
(menjaga privasi; warisi Paper~4). Dataset CSE-CIC-IDS2018 \cite{cicids} dan
UNSW-NB15 \cite{unsw} dipartisi menjadi $K{=}10$ klien melalui partisi non-IID
berbasis Dirichlet; sebagian klien ditetapkan sebagai penyerang.

\subsection{Model ancaman}
\label{sec:threat}
Kami mengasumsikan \emph{honest-but-curious} server dan hingga $f$ klien jahat
($f \in \{0,1,2,3\}$ dari $K{=}10$, yakni 0--30\%). Tiga serangan:
\begin{itemize}
  \item \textbf{\emph{Label-flipping}:} klien jahat membalik label latih lokalnya
  (attack$\leftrightarrow$benign) sebelum pelatihan.
  \item \textbf{Penskalaan bobot:} klien jahat mengirim pembaruan yang diperbesar
  faktor $\gamma>1$ untuk mendominasi rata-rata.
  \item \textbf{\emph{Backdoor}:} klien jahat menyisipkan pola pemicu (\emph{trigger})
  pada subset fitur dan memaksa label target (attack$\rightarrow$benign), diukur
  dengan \emph{Attack Success Rate} (ASR).
\end{itemize}

\subsection{Aturan agregasi yang dibandingkan}
\label{sec:agg}
\begin{itemize}
  \item \textbf{FedAvg} \eqref{eq:fedavg}: rata-rata berbobot (baseline rentan).
  \item \textbf{Krum / Multi-Krum} \cite{fkrum}: pilih pembaruan yang paling dekat
  ke mayoritas (jumlah jarak kuadrat ke $K{-}f{-}2$ tetangga terdekat minimum).
  \item \textbf{\emph{Trimmed-Mean}} \cite{ftrimmed}: per-koordinat, buang $\beta$
  nilai terkecil dan terbesar, rata-ratakan sisanya.
  \item \textbf{\emph{Median}} per-koordinat \cite{ftrimmed}.
\end{itemize}

\begin{align}
  w_{t+1}^{\text{FedAvg}} &= \sum_{k} \tfrac{n_k}{n}\, w_k^{(t)} \label{eq:fedavg}\\[2pt]
  \mathrm{Krum} &= \arg\min_{i}\ \textstyle\sum_{j \in \mathcal{N}_i} \lVert w_i - w_j \rVert^2 \label{eq:krum}
\end{align}

\noindent dengan $\mathcal{N}_i$ himpunan $K{-}f{-}2$ tetangga terdekat dari
pembaruan~$i$. Metrik utama MCC \eqref{eq:mcc}:
\begin{equation}
  \mathrm{MCC} = \frac{TP\cdot TN - FP\cdot FN}
  {\sqrt{(TP{+}FP)(TP{+}FN)(TN{+}FP)(TN{+}FN)}}, \label{eq:mcc}
\end{equation}
dan untuk \emph{backdoor}, $\mathrm{ASR} = N_{\text{trigger}\rightarrow\text{target}} / N_{\text{trigger}}$.

\subsection{Protokol evaluasi}
\label{sec:protocol}
Lima \emph{seed} $\{13,42,101,202,303\}$; rerata $\pm$ simpangan baku dan selang
kepercayaan 95\% (pola Paper~2/4). Untuk tiap kombinasi (serangan $\times$ fraksi
$\times$ aturan agregasi) dicatat MCC GLOBAL-\emph{test}, \emph{recall}/\emph{precision}/F1,
\emph{balanced accuracy}, dan ASR (khusus \emph{backdoor}). Logika agregasi
diimplementasikan di atas \texttt{aws/fed\_node.py} agar simulasi konsisten dengan
penggelaran EC2.

% ============================================================
\FloatBarrier
\section{Hasil}
\label{sec:results}

\begin{quote}
\small\itshape Catatan kejujuran: seluruh \emph{tabel} hasil (Tabel~\ref{tab:vuln}--\ref{tab:clean})
berisi sel \TBD{} dan akan diisi HANYA dari eksperimen nyata (notebook
\texttt{42\_01--42\_07} + eksekusi AWS). Subbagian \ref{sec:expected} di bawah memuat
\textbf{hipotesis kuantitatif ilustratif} --- angka yang kami \emph{harapkan}, bukan
hasil pengukuran --- semata untuk menjelaskan arah dan bentuk kontribusi. Angka
ilustratif TIDAK muncul di tabel hasil dan akan dikonfirmasi/dikoreksi oleh data.
\end{quote}

\subsection{Hipotesis kuantitatif (ILUSTRATIF --- bukan hasil pengukuran)}
\label{sec:expected}

\begin{quote}
\small\itshape Peringatan: seluruh angka pada subbagian ini adalah \textbf{ekspektasi
hipotetis} untuk memberi gambaran arah paper. Bukan hasil eksperimen. Nilai nyata
akan menggantikannya setelah run.
\end{quote}

\noindent Berdasarkan rancangan dan temuan Paper~4 (FedAvg lintas-\emph{dataset}
mencapai MCC sekitar $0{,}65$ pada kondisi jujur), kami merumuskan pola yang
\emph{diharapkan}:

\begin{itemize}
  \item \textbf{RQ1 (kerentanan FedAvg).} Kami \emph{menduga} MCC FedAvg menurun
  monoton seiring naiknya fraksi penyerang: dari acuan $\sim\!0{,}65$ ($f{=}0$)
  menjadi $\sim\!0{,}45$ ($f{=}10\%$), $\sim\!0{,}25$ ($f{=}20\%$), hingga $<\!0{,}15$
  ($f{=}30\%$). Penskalaan bobot diduga paling merusak, \emph{label-flipping} paling ringan.
  \item \textbf{RQ2 (pemulihan oleh robust aggregation).} Pada $f{=}20\%$, kami
  \emph{menduga} Krum/Multi-Krum/\emph{Trimmed-Mean}/\emph{Median} mempertahankan MCC
  di kisaran $\sim\!0{,}55$--$0{,}62$ (dekat kondisi jujur), jauh di atas FedAvg
  ($\sim\!0{,}25$). Untuk \emph{backdoor}, aturan tangguh diharapkan menekan ASR dari
  $\sim\!80$--$95\%$ (FedAvg) menjadi $<\!10\%$.
  \item \textbf{RQ3 (harga ketahanan).} Pada $f{=}0$, kami \emph{menduga} aturan
  tangguh membayar \emph{clean cost} kecil: MCC turun $\sim\!0{,}02$--$0{,}05$
  dibanding FedAvg (mis.\ $0{,}65 \rightarrow 0{,}60$--$0{,}63$). Krum diduga membayar
  paling mahal (membuang satu klien penuh), \emph{Trimmed-Mean}/\emph{Median} paling murah.
\end{itemize}

\noindent Pola di atas adalah \emph{kriteria keberhasilan} yang akan diuji: bila data
nyata mendukungnya, kontribusi paper menguat; bila tidak (mis.\ heterogenitas
lintas-\emph{dataset} membuat aturan tangguh keliru membuang klien jujur), temuan
tersebut dilaporkan apa adanya sebagai kontribusi tersendiri.

\subsection{Kerentanan FedAvg terhadap fraksi penyerang (RQ1)}
Tabel~\ref{tab:vuln} menunjukkan MCC GLOBAL-\emph{test} FedAvg seiring naiknya
fraksi penyerang untuk tiap jenis serangan. Diharapkan MCC menurun saat fraksi naik.

\begin{table}[t]
\centering
\caption{MCC FedAvg vs fraksi penyerang (GLOBAL-\emph{test}, rerata 5 \emph{seed}).
Dari \texttt{42\_04}.}
\label{tab:vuln}
\begin{tabular}{lcccc}
\toprule
Serangan & $f{=}0\%$ & $f{=}10\%$ & $f{=}20\%$ & $f{=}30\%$ \\
\midrule
\emph{Label-flipping} & \TBD & \TBD & \TBD & \TBD \\
Penskalaan bobot      & \TBD & \TBD & \TBD & \TBD \\
\emph{Backdoor}       & \TBD & \TBD & \TBD & \TBD \\
\bottomrule
\end{tabular}
\end{table}

\subsection{Pemulihan oleh robust aggregation (RQ2)}
Tabel~\ref{tab:defense} membandingkan MCC lima aturan agregasi pada fraksi penyerang
tetap (mis.\ 20\%) untuk tiap serangan. Diharapkan aturan tangguh mempertahankan MCC
jauh di atas FedAvg.

\begin{table}[t]
\centering
\caption{MCC per aturan agregasi pada $f{=}20\%$ (GLOBAL-\emph{test}). Dari \texttt{42\_05}.}
\label{tab:defense}
\begin{tabular}{lccc}
\toprule
Agregasi & \emph{Label-flip} & Penskalaan & \emph{Backdoor} \\
\midrule
FedAvg (rentan) & \TBD & \TBD & \TBD \\
Krum            & \TBD & \TBD & \TBD \\
Multi-Krum      & \TBD & \TBD & \TBD \\
\emph{Trimmed-Mean} & \TBD & \TBD & \TBD \\
\emph{Median}   & \TBD & \TBD & \TBD \\
\bottomrule
\end{tabular}
\end{table}

\subsection{Efektivitas backdoor: ASR (RQ2)}
Tabel~\ref{tab:asr} melaporkan \emph{Attack Success Rate} \emph{backdoor}. Aturan
tangguh diharapkan menekan ASR mendekati nol, sedangkan FedAvg tinggi.

\begin{table}[t]
\centering
\caption{ASR \emph{backdoor} per aturan agregasi vs fraksi penyerang. Dari \texttt{42\_06}.}
\label{tab:asr}
\begin{tabular}{lccc}
\toprule
Agregasi & $f{=}10\%$ & $f{=}20\%$ & $f{=}30\%$ \\
\midrule
FedAvg      & \TBD & \TBD & \TBD \\
Krum        & \TBD & \TBD & \TBD \\
\emph{Trimmed-Mean} & \TBD & \TBD & \TBD \\
\emph{Median} & \TBD & \TBD & \TBD \\
\bottomrule
\end{tabular}
\end{table}

\subsection{Harga ketahanan tanpa serangan (RQ3)}
Tabel~\ref{tab:clean} membandingkan MCC tiap aturan saat $f{=}0\%$ (tanpa penyerang).
Selisih terhadap FedAvg = \emph{clean cost} (harga ketahanan).

\begin{table}[t]
\centering
\caption{MCC saat $f{=}0\%$ (harga ketahanan). Dari \texttt{42\_05}/\texttt{42\_07}.}
\label{tab:clean}
\begin{tabular}{lcc}
\toprule
Agregasi & MCC GLOBAL ($f{=}0$) & $\Delta$ vs FedAvg \\
\midrule
FedAvg      & \TBD & --- \\
Krum        & \TBD & \TBD \\
Multi-Krum  & \TBD & \TBD \\
\emph{Trimmed-Mean} & \TBD & \TBD \\
\emph{Median} & \TBD & \TBD \\
\bottomrule
\end{tabular}
\end{table}

% ============================================================
\FloatBarrier
\section{Diskusi}
\label{sec:discussion}

Jika hipotesis terkonfirmasi, kontribusi utama bukanlah ``robust aggregation
bekerja'' --- itu diketahui pada data IID --- melainkan \emph{karakterisasi}
perilakunya pada FL-NIDS lintas-\emph{dataset} yang sudah menantang bahkan tanpa
penyerang (temuan Paper~4). Pertanyaan menarik: apakah heterogenitas lintas-\emph{dataset}
memperburuk kerentanan (penyerang lebih mudah bersembunyi di antara variasi wajar),
dan apakah aturan tangguh yang membuang \emph{outlier} justru keliru membuang klien
jujur yang distribusinya ekstrem. Hasil nyata akan menjawabnya; kami laporkan apa
adanya, termasuk bila ada aturan yang gagal pada setting ini.

% ============================================================
\FloatBarrier
\section{Keterbatasan}
\label{sec:limitations}

Pertama, model ancaman terbatas pada tiga serangan representatif; serangan adaptif
yang mengetahui aturan agregasi tidak dicakup. Kedua, eksperimen memakai dua
\emph{dataset}/sepuluh klien sebagai instansiasi konkret; generalisasi ke lebih
banyak klien/alat adalah kerja lanjutan. Ketiga, aturan berbasis \emph{trust}
(mis.\ FLTrust) yang memerlukan \emph{root-set} tepercaya di server tidak menjadi
fokus utama. Keempat, seluruh angka akan berasal dari eksperimen nyata; sebelum itu
tidak ada klaim kuantitatif.

% ============================================================
\FloatBarrier
\section{Kesimpulan dan Kerja Mendatang}
\label{sec:conclusion}

Makalah ini merancang evaluasi \emph{robust aggregation} untuk FL-NIDS
lintas-\emph{dataset} di ruang SFM, melawan \emph{poisoning} dan \emph{backdoor}.
Draf ini memaparkan model ancaman, aturan agregasi, dan protokol evaluasi; seluruh
angka hasil akan diisi dari eksperimen nyata sebelum publikasi. Kerja mendatang
meliputi serangan adaptif, aturan berbasis \emph{trust}, dan kombinasi dengan
personalisasi (temuan mitigasi Paper~4).

% ============================================================
\section*{Ketersediaan Kode dan Data}
Implementasi node federasi dan aturan agregasi (\texttt{aws/fed\_node.py} +
\texttt{notebook42/}), serta \emph{notebook} eksperimen (\texttt{42\_01--42\_07})
disertakan pada repositori proyek. Ruang SFM diwarisi dari Paper~1. Dataset bersifat
publik \cite{cicids,unsw}.

\section*{Pernyataan Etika}
Seluruh serangan disimulasikan dalam lingkungan riset terkontrol memakai
\emph{dataset} publik, bertujuan memperkuat pertahanan NIDS. Tidak ada data pribadi
yang dikumpulkan dan tidak ada sistem produksi yang diserang.

% ============================================================
\begin{thebibliography}{99}
% Format IEEE. Nomor urut sesuai kemunculan.

\bibitem{sfm} H.~Y.~Martono, I.~Syarif, and F.~A.~Saputra, ``Cross-dataset
generalization of ML-based network intrusion detection via semantic feature
mapping and few-shot adaptation, validated on real cloud traffic,'' \emph{status
terbit diisi saat accepted}.

\bibitem{popoola} S.~I.~Popoola, G.~Gui, B.~Adebisi, M.~Hammoudeh, and
H.~Gacanin, ``Federated deep learning for collaborative intrusion detection in
heterogeneous networks,'' in \emph{Proc. IEEE 94th Veh. Technol. Conf.
(VTC-Fall)}, 2021, pp.~1--6.

\bibitem{fsurvey1} A.~Khraisat, A.~Alazab, T.~Jan, and A.~Gopez, ``Survey on
federated learning for intrusion detection system,'' \emph{ACM Comput. Surv.},
2024.

\bibitem{fsurvey2} J.~L.~Hern\'andez-Ramos \emph{et~al.}, ``Intrusion detection
based on federated learning: A systematic review,'' \emph{ACM Comput. Surv.}, 2025.

\bibitem{ftaxonomy} A.~A.~Wardana and P.~Sukarno, ``Taxonomy and survey of
collaborative intrusion detection system using federated learning,'' \emph{ACM
Comput. Surv.}, 2024.

\bibitem{fbackdoor} Z.~Sun, P.~Kairouz, A.~T.~Suresh, and H.~B.~McMahan, ``Can you
really backdoor federated learning?'' \emph{arXiv:1911.07963}, 2019.

\bibitem{fbagdasaryan} E.~Bagdasaryan, A.~Veit, Y.~Hua, D.~Estrin, and
V.~Shmatikov, ``How to backdoor federated learning,'' in \emph{Proc. AISTATS},
2020, pp.~2938--2948.

\bibitem{ffung} C.~Fung, C.~J.~M.~Yoon, and I.~Beschastnikh, ``Mitigating sybils
in federated learning poisoning,'' \emph{arXiv:1808.04866}, 2018.

\bibitem{fkrum} P.~Blanchard, E.~M.~El~Mhamdi, R.~Guerraoui, and J.~Stainer,
``Machine learning with adversaries: Byzantine tolerant gradient descent,'' in
\emph{Proc. NeurIPS}, 2017, pp.~119--129.

\bibitem{ftrimmed} D.~Yin, Y.~Chen, R.~Kannan, and P.~Bartlett, ``Byzantine-robust
distributed learning: Towards optimal statistical rates,'' in \emph{Proc. ICML},
2018, pp.~5650--5659.

\bibitem{sarhan} M.~Sarhan, S.~Layeghy, and M.~Portmann, ``Towards a standard
feature set for network intrusion detection system datasets,'' \emph{Mobile Netw.
Appl.}, vol.~27, no.~1, pp.~357--370, 2022.

\bibitem{cicids} I.~Sharafaldin, A.~H.~Lashkari, and A.~A.~Ghorbani, ``Toward
generating a new intrusion detection dataset and intrusion traffic
characterization,'' in \emph{Proc. 4th ICISSP}, 2018, pp.~108--116.

\bibitem{unsw} N.~Moustafa and J.~Slay, ``UNSW-NB15: A comprehensive data set for
network intrusion detection systems,'' in \emph{Proc. MilCIS}, 2015, pp.~1--6.

\end{thebibliography}

\end{document}
